\SourceLeaf{71}
\[
\cos\varphi^r=(1+a)\frac{1}{2^r}.{}^{r}\mathfrak{B}^{(\frac12(r+a))}\cos a\varphi+.\,.\,.+\frac{1}{2^{r-1}}\mathfrak{B}^{(\frac12(r+a)+h)}\cos(2h+a)\varphi+.\,.\,.,
\]

\noindent atque sic binae series, quae primo adspectu diversissimae videbantur, sub una expressione, eaque satis simplici, comprehendi possunt.

\Needspace{10\baselineskip}\par\smallskip\hrulefill\smallskip\par

\begin{center}\textbf{BEMERKUNGEN.}\end{center}

Das Manuscript der vorliegenden »Dissertatiuncula« wurde 1893 von Herrn Prof. \textsc{Wilhelm Meyer} (Göttingen) in den Acten der Königl. Gesellschaft der Wissenschaften zu Göttingen gefunden und ist von E. \textsc{Schering} mit einer Reihe erläuternder Bemerkungen in den »Göttinger Nachrichten« von 1893 No.\,15 veröffentlicht. Es fehlt in der Urschrift der Name des Verfassers und ebenso auch eine Angabe über Ort und Zeit der Abfassung. Indessen beseitigte der Vergleich sonstiger Manuscripte von \textsc{Gauss} mit der vorliegenden Handschrift jeden Zweifel, welcher sonst vielleicht an \textsc{Gauss}' Autorschaft bestehen könnte.

Sachlich ist zu bemerken, dass \textsc{Gauss} vermuthlich schon zu der Zeit, als er das Collegium Carolinum zu Braunschweig besuchte (1792 bis 1795), die Werke von \textsc{Lagrange} und \textsc{Euler} studirte. Es finden sich im Nachlass (wahrscheinlich aus dieser Zeit stammend) verschiedene Auszüge aus Werken von \textsc{Lagrange} und \textsc{Euler}, und darunter insbesondere solche aus \textsc{Euler}s Institutiones calculi integralis, Theil I Kapitel 6, an welche die Dissertatiuncula unmittelbar anknüpft.

Ist hiernach auch sachlich \textsc{Gauss}' Verfasserschaft unzweifelhaft, so wird man gleichwohl die fast vollständige Abhängigkeit der ganzen Entwicklung von \textsc{Euler} nur mit der Annahme vereinigen können, dass die Entstehung der fraglichen Untersuchung entweder in den ersten Anfang der Göttinger Studienzeit von \textsc{Gauss} (Herbst 1795 bis 1798) oder auch noch vor diese Zeit fällt. Es löst nämlich \textsc{Euler} die Aufgabe der Entwicklung von $(1+n\cos\varphi)^\nu$ in die Reihe

\[
A + A^{(1)}\cos\varphi + A^{(2)}\cos 2\varphi + A^{(3)}\cos 3\varphi + .\,.\,.
\]

\noindent so, dass er nur $A$ und $A^{(1)}$ direct angiebt, sich aber wegen der weiter folgenden Coefficienten auf eine Recursionsformel beruft, welche durch einen naheliegenden Kunstgriff gewonnen wird. Unter Vermeidung des letzteren giebt \textsc{Gauss} auch die höheren Coefficienten $A^{(2)}$,\,.\,. durch directe Betrachtung, reproducirt aber übrigens alle Ansätze und Weiterentwicklungen \textsc{Euler}s.

Die Citate auf \textsc{Kästner} sprechen dafür, dass die definitive Abfassung der Dissertatiuncula in den Beginn der Göttinger Studienzeit fällt. Die gegen Ende von No.\,10 genannte Abhandlung \textsc{Klügel}s gehört den mathematischen Commentationen der Göttinger Societät von 1793 und 94 an. Dieselben sind als Gesammtband allerdings erst 1796 ausgegeben; doch sind \textsc{Gauss} die einzelnen Abhandlungen wahrscheinlich schon früher zugänglich gewesen.

Für die Annahme, dass die Entstehung der Dissertatiuncula spätestens Anfang 1796 anzusetzen ist, würde auch der Umstand sprechen, dass in dem Ende März 1796 beginnenden wissenschaftlichen Tagebuche keine Notiz über diese Abhandlung gefunden wird.

\Needspace{10\baselineskip}
Eine Einzelbemerkung erfordert noch die Entwicklung unter Nr.\,8, $\gamma$. Der Vorwurf eines Fehlers bei \textsc{Euler} ist nicht berechtigt; vielmehr hat sich \textsc{Gauss} versehen, indem er bei der Berechnung der Gleichungen für $A^{(h)}$, $A^{(h+1)}$ rechter Hand den Factor $\frac{2}{\BareRoot{(1-nn)}}$ übersah. Indem \textsc{Schering} in der von ihm veranstalteten Ausgabe diesen Factor zufügte, stellte er die Übereinstimmung mit \textsc{Euler} her. Der vorliegende Abdruck gibt gleichfalls die berichtigten Formeln.\hfill\textsc{Fricke}.

\par\smallskip\hrulefill\smallskip\par

